\begin{proofbox}
	\textbf{\textcolor{CProof}{思路提示}}

	采用反证法：假设 $a$ 与 $l$ 不平行，利用它们在同一个平面 $\beta$ 内必相交，导出与线面平行矛盾。

	\medskip
	\textbf{\textcolor{CProof}{正式推导}}
	\begin{description}[style=nextline, leftmargin=2.8em, labelsep=0.5em, itemsep=0.55em]
		\item[\textbf{步骤一（反设结论）：}] 假设 $a$ 与 $l$ 不平行.
			\[
				a\nparallel l.
			\]
			\par\textbf{理由：}
			在平面 $\beta$ 内两直线要么平行要么相交，若要证明平行，先假设不平行。

		\item[\textbf{步骤二（相交得公共点）：}] 由 $a,l\subset\beta$ 且不平行，故 $a\cap l$ 存在，设 $P=a\cap l$.
			\[
				P=a\cap l.
			\]
			\par\textbf{理由：}
			平面内不平行则相交，交点唯一。

		\item[\textbf{步骤三（导出与已知矛盾）：}] 因为 $l\subset\alpha$，所以 $P\in\alpha$，又 $P\in a$，故 $P\in a\cap\alpha$.
			\[
				P\in a,\; P\in l\subset\alpha \Rightarrow P\in a\cap\alpha.
			\]
			\par\textbf{理由：}
			由 $l\subset\alpha$ 得 $P\in\alpha$，从而 $a$ 与 $\alpha$ 有公共点 $P$，但已知 $a\parallel\alpha$ 意味着 $a\cap\alpha=\varnothing$，矛盾。

		\item[\textbf{步骤四（结论成立）：}] 因此反设错误，必有 $a\parallel l$.
			\[
				\therefore a\parallel l.
			\]
			\par\textbf{理由：}
			反证法，假设不成立，故 $a\parallel l$.
	\end{description}

	\medskip
	\textbf{\textcolor{CProof}{结论回扣}}

	由上述反证法严格证明了线面平行性质定理：若直线 $a$ 平行于平面 $\alpha$，且 $a$ 在平面 $\beta$ 内，$\alpha$ 与 $\beta$ 的交线为 $l$，则 $a\parallel l$.
\end{proofbox}
